\documentclass[10pt]{article}
\usepackage[margin=1in]{geometry}
\usepackage{amsmath,amssymb,amsthm}
\usepackage[hidelinks]{hyperref}
\newtheorem{theorem}{Theorem}
\newtheorem{corollary}[theorem]{Corollary}
\newcommand{\ds}{\,ds}
\newcommand{\NN}{\mathbb R^2}
\begin{document}
\begin{center}
{\Large Circular asymptotics of eventually convex planar free elastic flows}
\end{center}
\begin{abstract}
For a smooth free elastic flow of closed plane curves, the product of
enclosed area and squared bending energy decreases strictly at each
embedded convex time slice unless that slice is round. This classifies
convex similarity solutions and gives smooth convergence, after length
normalization, for every flow which is embedded and convex for all
sufficiently large times. The proof combines Gage's inequality with an
exact dissipation identity and the circle stability theorem of Miura and
Wheeler.
\end{abstract}

\section{Area and bending energy}
Orient an embedded convex time slice counterclockwise, and let \(s\) be
arclength, \(k\geq 0\) its curvature and \(N\) its inward unit normal.
The free elastic flow has normal velocity \(-F N\), where
\[
 F=k_{ss}+\tfrac12 k^3,\qquad E=\tfrac12\int k^2\ds .
\]
Write \(A\) for the enclosed area, \(L\) for the length,
\(B=\int k^2\ds=2E\), \(C=\int k^3\ds\), \(D=\int F^2\ds\), and
\(Q=AE^2\). We integrate around the entire closed curve. At a convex
slice \(C>0\) since \(\int k\ds=2\pi\). The area variation and the
gradient-flow identity give
\[
 A'=\int F\ds=C/2,\qquad E'=-D.
\]
These calculations are independent of tangential velocity.

\begin{theorem}\label{mon}
At every embedded convex time slice \(Q'\leq0\), with equality precisely
on circles. On every time interval of embedded convexity, set
\(\tau=\log(A/A_*)\) with \(A_*>0\) fixed. Then
\begin{equation}\label{factor}
 -\frac{d\log Q}{d\tau}=
 \frac{AB}{\pi L}\ \frac{2\pi C}{B^2}\ \frac{4LD}{C^2}-1 .
\end{equation}
Each of the three factors is at least one.
\end{theorem}
\begin{proof}
The first variations give the left side as \(4AD/(EC)-1\), which is
the right side of \eqref{factor}. Gage's inequality \(AB\geq\pi L\)
holds for closed \(C^2\) convex boundaries \cite{gage}. Cauchy--Schwarz
applied to \(k^{1/2}\) and \(k^{3/2}\) gives \(B^2\leq2\pi C\);
applied to \(1\) and \(F\), it gives \(C^2\leq4LD\).
If equality holds in \eqref{factor} in the sense that \(Q'=0\), then
\(k^{3/2}=\lambda k^{1/2}\) for some constant \(\lambda>0\).
The continuous function \(k\) takes its values in \(\{0,\lambda\}\).
Connectedness and total curvature \(2\pi\) give \(k\equiv\lambda\).
The three factors equal one on a circle.
\end{proof}

\begin{corollary}\label{breath}
If two time slices of a free elastic flow are related by a Euclidean
similarity and every intermediate slice is an embedded convex curve,
the flow on that interval is a family of concentric expanding circles.
In particular, every embedded convex similarity solution is round.
\end{corollary}
\begin{proof}
The similarity invariance of \(Q\) gives equal endpoint values.
Theorem~\ref{mon} forces every intermediate slice to be round.
For circles of center \(c(t)\), radius \(r(t)\) and outward unit normal
\(n\), the normal equation reads
\(\langle c'(t),n\rangle+r'(t)=1/(2r(t)^3)\) in every direction \(n\).
Comparing opposite directions gives \(c'=0\), then
\((r^4)'=2\). In particular all time slices have a common point in their
convex hulls. The same conclusion holds for covers of such a boundary,
by working on the curve traced once.
\end{proof}

\section{Convergence from eventual convexity}
\begin{theorem}\label{conv}
Let \(\gamma\) be a smooth closed planar free elastic flow. Suppose there
is a time \(t_0\) such that every slice at \(t\geq t_0\) is an embedded
convex curve. Then \(\gamma(\cdot,t)/L(t)\) converges smoothly up to
reparametrization to the circle of radius \(1/(2\pi)\) centered at the
origin. Furthermore
\[
 \lim_{t\to\infty}\frac{L(t)^4}{t}=32\pi^4,\qquad
 \lim_{t\to\infty}t^{1/4}E(t)=\frac{\pi}{2^{1/4}} .
\]
\end{theorem}
\begin{proof}
Gage's inequality and the isoperimetric inequality \(L^2\geq4\pi A\)
imply
\[
 \pi^3\leq Q(t)\leq Q(t_0),\qquad
 LB\leq\frac{4Q(t_0)}{\pi},\qquad
 C\geq\frac{B^2}{2\pi}\geq\frac{\pi L^2}{2A^2}
      \geq\frac{2\pi^2}{A}.
\]
Thus \((A^2)'=AC\geq2\pi^2\), and
\(\tau=\log(A/A(t_0))\) increases to infinity.
The nonnegative expression in \eqref{factor} has finite integral in
\(\tau\), bounded by \(\log(Q(t_0)/\pi^3)\).
Choose \(\tau_j\to\infty\) so the expression tends to zero.
Each of its three factors then tends to one.

Rescale each selected slice to length one and use unit-speed coordinates
on \(\mathbb R/\mathbb Z\). The three factors are scale invariant. The
rescaled \(B_j\) equals the original \(L B\), hence stays bounded.
The second and third factors in turn bound the rescaled \(C_j\) and
\(D_j\). Thus \(\|k_j\|_2\) and \(\|F_j\|_2\) are uniformly bounded.
To pass to a compact limit, put \(X_j=\|k_j''\|_2\).
Integration by parts gives \(\|k_j'\|_2^2\leq \|k_j\|_2 X_j\).
Some point of the unit period satisfies \(k_j^2\leq \|k_j\|_2^2\);
integrating \((k_j^2)'=2k_j k_j'\) from that point gives
\[
 \|k_j\|_\infty^2\leq \|k_j\|_2^2+2\|k_j\|_2\|k_j'\|_2
 \leq c(1+\sqrt{X_j}).
\]
The equation \(k_j''=F_j-k_j^3/2\) now yields
\[
 X_j\leq \|F_j\|_2+\tfrac12\|k_j\|_\infty^2\|k_j\|_2
 \leq c'(1+\sqrt{X_j}).
\]
This bounds \(X_j\). The resulting uniform \(H^2\) bound gives a
subsequence converging in \(C^1\) to a nonnegative periodic function
\(k_*\) of integral \(2\pi\). Passing to the limit in the second factor
gives \(2\pi\int k_*^3=(\int k_*^2)^2\). As in Theorem~\ref{mon}, equality
forces \(k_*\equiv 2\pi\).

At the corresponding original times \(t_j\) we therefore have
\[
 L(t_j)^3\int_{\gamma(t_j)} k_s^2\ds
   =\int_0^1 |k_j'|^2\,ds \longrightarrow 0 .
\]
Restart the flow at one of these times below the turning-number-one
threshold in \cite[Theorem 1.2]{mw}. That theorem gives the stated smooth
convergence, including the center of the normalized limit.
The length variation is \(L'=\int kF\ds=-\int k_s^2\ds+\tfrac12\int
k^4\ds\). Smooth convergence in length-one variables yields
\(L^3L'\to(2\pi)^4/2\) and \(LE\to2\pi^2\).
Integrating \((L^4)'=4L^3L'\) proves the two limits.
\end{proof}

For any fixed initial slice, its image lies in the convex hull of every
sufficiently late slice under the hypotheses of Theorem~\ref{conv}.
Indeed the normalized limit circle has support \(1/(2\pi)\) in every
direction. Uniform convergence therefore puts the disk of radius
\(L(t)/(4\pi)\) about the origin in the unscaled convex hull for all
sufficiently large \(t\). This conclusion leaves possible crossings
between earlier slices open. Eventual convexity is a hypothesis on the
trajectory. The unrestricted immersed flow in Section 1.2 of \cite{mw}
can have inflections and self-intersections.

\begin{thebibliography}{9}
\bibitem{gage} M. E. Gage,
An isoperimetric inequality with applications to curve shortening,
\emph{Duke Math. J.} 50 (1983), 1225--1229.
\href{https://doi.org/10.1215/S0012-7094-83-05052-4}{doi:10.1215/S0012-7094-83-05052-4}.
\bibitem{mw} T. Miura and G. Wheeler,
The free elastic flow for closed planar curves,
\emph{J. Funct. Anal.} 289 (2025), article 111030.
\href{https://arxiv.org/html/2404.12619v3}{arXiv:2404.12619v3}.
\end{thebibliography}
\end{document}
